\chapter{The series}

\textbf{Purpose:} Report the controlled series and every number in it, including the rows the
detector cannot be asked to read. \textbf{Scope:}
\texttt{examples/\allowbreak{}game\_\allowbreak{}theory/\allowbreak{}4\_\allowbreak{}measure/\allowbreak{}grundy\_\allowbreak{}period.\allowbreak{}py}

\section{How a row is made}

A move set is a finite set of positive integers. The game is one heap, a move takes a number in that
set, and whoever cannot move loses. Two thousand and forty eight positions are generated, each
carrying the least non-negative integer absent from the values it reaches.

The answer is computed twice by two routines that share nothing. The exact one slides the tail of
the sequence against itself and returns the shortest shift under which it repeats to the end. The
measured one is \texttt{measure.\allowbreak{}periodicity.\allowbreak{}sequence\_\allowbreak{}period},
the same function the crystallography chapter asks for a cell edge, which is handed a list of small
integers and told nothing else.

The sweep is every move set of one to four moves drawn from 1 to 10, which is three hundred and
eighty five games. Every one of them settled: the exact routine found a period on all three hundred
and eighty five, and the periods present run 2 to 45.

\section{The window is a constraint and not a result}

The detector scores a candidate period together with its own multiples, which is the correction the
crystallography chapter had to make when published cell edges came back at twice their value on ten
of eighteen axes. A period beyond half the window has one multiple inside the window, so its family
is one lag and the score degenerates to taking the tallest lag again.

Rows whose true period exceeds half the window are therefore out of range by construction. They are
counted separately below and they are not failures. Counting them as failures would be reporting the
detector's documented precondition as its error rate.

\begin{center}
\begin{tabular}{@{}rrrr@{}}
\toprule
window & in range & exact & out of range \\
\midrule
16 & 115 & 115 & 270 \\
24 & 220 & 220 & 165 \\
32 & 346 & 346 & 39 \\
48 & 378 & 378 & 7 \\
64 & 383 & 383 & 2 \\
\bottomrule
\end{tabular}
\end{center}

Three hundred and eighty three of three hundred and eighty three at the widest window, and the same
at every narrower one over the rows it can reach. The two rows still out of range at a window of 64
have periods of 45 and 33.

\section{Every row against its own floor}

A count of agreements is not evidence on its own, because a detector that always answered 2 would
agree with every period-2 row. Each sequence is therefore scored against twelve shuffles of itself.
A shuffle keeps every Grundy value and destroys which position carries it, so whatever margin it
reaches is what the detector reads off the histogram alone, and the figure quoted is the tallest of
the twelve rather than their average.

At a window of 64, over three hundred and eighty three rows:

\begin{center}
\begin{tabular}{@{}lrrr@{}}
\toprule
 & minimum & median & maximum \\
\midrule
live margin & 1.016 & 1.091 & 2.000 \\
shuffle floor & 0.030 & 0.044 & 0.066 \\
\bottomrule
\end{tabular}
\end{center}

The worst ratio of live margin to its own floor is 16.0 and the median is 25.6. No row failed to
clear its floor. The bands do not touch: the smallest live margin is fifteen times the largest
floor anywhere in the series.

\section{What the series does and does not license}

It licenses a statement about this domain. Over subtraction games with move sets drawn from 1 to 10,
the detector returns the exact Grundy period on every sequence whose period it is able to score. The
quantifier is over an enumerated set, the answers came from a proof, and the count is large enough
that a coincidence would have to be systematic rather than lucky.

It does not license a statement about any other domain, and it does not upgrade the workbook's span
table. The workbook's own deduction says why: soundness belongs to the construction and cost belongs
to the domain's distribution, so a result obtained inside one distribution stays inside it. A
subtraction game's Grundy sequence is exactly periodic after a short preperiod, with two to a few
dozen distinct values and no noise at all. That is the easiest thing a period detector will ever be
shown. The series establishes that the detector is right where an answer exists and is checkable,
at a count nobody has to take on faith, and the next chapter is the half that matters more.
